\documentclass[../main.tex]{subfiles}

\begin{document}
\begin{problema}
    Sean \(f,g \colon \Omega \subseteq \mathbb{C} \to \mathbb{C}\) analíticas sobre la
    región \(\Omega\) y continuas sobre \(\overline{\Omega}\) tal que \(\Omega\) es
    compacto, tales que \(f(z) = g(z)\) para toda \(z\in \mathrm{Fr}(\Omega)\). Pruebe
    que \(f(z) = g(z)\) para toda \(z\in \Omega\).
\end{problema}
\startsolution

Vamos a usar el corolario del módulo máximo aplicado a la función \(h = f - g\). Notemos
que \(\overline{\Omega}\) es compacto, por lo que \(\Omega\) es una región acotada,
\(h\) es continua en \(\overline{\Omega}\) y analítica en \(\Omega\). Por lo que se
satisfacen las hipótesis del corolario. Es decir, \(\abs{h}\) alcanza su máximo en
\(\partial\Omega\).

Pero \( h \equiv 0\) en \(\partial\Omega\) pues por hipótesis, \(f(z) = g(z)\quad
\forall z\in \partial\Omega\). Entonces, \(\max_{\partial\Omega}\abs{h} = 0\) y

\begin{equation*}
    \sup_{\Omega}\abs{h} \leq \max_{\overline{\Omega}}\abs{h} =
    \max_{\Omega}\abs{h} = 0.
\end{equation*}

Esto implica que \(\sup_{\Omega}\abs{h} = 0\) y además que \(h(z) = 0,\; \forall z\in
\Omega\). Por lo tanto,

\begin{equation*}
    f(z) = g(z)\quad \forall z\in \Omega.
\end{equation*}
\end{document}
